% 原文 pp.300–306（印刷 pp.1418–1424）
% 本文件承接 sol_p1.tex。p.300 起为「习题 10 的解答」的后半段，该页中部转入「习题 11 的解答」。

% ===== 原文 p.300 =====
\subsection{习题 10 的解答（续）}

如问题陈述中所说明，我们只取单个可观测量，它是 $P(\phi)$ 的一个多项式。于是前述讨论表明，它的单点函数（one-point function）$E$ 归结为对 $\mathbb{C}^n$ 以及对奇变量（odd variables）的一个有限维积分：
\begin{equation}
  \tag{10.2}
  \label{eq:sol10-2}
  E=\int P(\phi)\,e^{-\|\operatorname{grad}W(\phi)\|^2+\frac12\operatorname{Hess}\overline{W}(\phi)(\bar\chi',\bar\chi)}
  \frac{d^n\bar\phi\,d^n\phi}{(2\pi i)^n}\,d^n\bar\chi\,d^n\bar\chi'.
\end{equation}

为使此式有意义，我们注意到 $d^n\bar\phi\,d^n\bar\chi$ 具有典范的意义，因为 $d^n\bar\phi$ 是 $\operatorname{Det}T^*_{0,1}$ 的一个截面，而 $d^n\bar\chi$ 是 $\operatorname{Det}T_{0,1}$ 的一个截面。然而 $d^n\phi$ 是 $\operatorname{Det}T^*_{1,0}$ 的截面，$d^n\bar\chi'$ 是 $\operatorname{Det}T_{0,1}\cong\operatorname{Det}T^*_{1,0}$ 的截面，其中最后一个同构用到了 $\mathbb{C}^n$ 上标准的 Hermite 度量。因此我们必须选定一个元素
\[
  \omega=d^n\phi\,d^n\bar\chi'\in(\operatorname{Det}T^*_{1,0})^{\otimes2},
\]
才能把积分 (10.2) 定义为一个数。

% [Tutor Note: (10.2) 的指数里 $\frac12\operatorname{Hess}\overline W(\phi)(\bar\chi',\bar\chi)$ 的两个 Grassmann 变量都带横杠、且其中一个带撇，撇不是求导而是区分两个独立变量的记号；若误把 $\bar\chi'$ 当作 $\bar\chi$ 的某种导数，就会把该二次型写成 $\operatorname{Hess}W$ 而非 $\operatorname{Hess}\overline W$，从而丢掉整个反全纯因子。另注意测度 $\frac{d^n\bar\phi\,d^n\phi}{(2\pi i)^n}$ 与 $d^n\bar\chi\,d^n\bar\chi'$ 的配对是 Berezin 积分已经做掉的部分，故下面「费米子积分很容易做出」只是把两个 Grassmann 变量积掉，留下 $\epsilon^{-n}\det\operatorname{Hess}\overline W$。]

% [Tutor Note: 末式 $E=\sum_{\text{crit pts}}\frac{P(\phi_i)}{\det\operatorname{Hess}W(\phi_i)}$ 并不是靠「Gauss 近似」凑出来的巧合，而是 Atiyah–Bott 不动点公式（等价于 Poincaré–Hopf 定理）在 $W$ 的梯度流上的特例：$\delta$-恰当的作用量使积分局域化，分母的行列式正是临界点的 Morse 指标符号。作者用「Gaussian 近似」的说法跳过了这一步，读者若想补足，可查 Atiyah–Bott, \emph{Topology} 6 (1967) 或 Witten 的《Supersymmetry and Morse theory》。]

如 §14.3.1 中所解释，拓扑关联函数（topological correlation function）与 lagrangian 中的绝大多数参数无关。的确，在这个有限维模型中 $L$ 是恰当的（exact），由 (10.1) 可以算出：
\[
  \bar\eta L=\delta\Bigl\{\frac{1}{2\sqrt2}\bigl\langle\bar\chi',\operatorname{grad}W(\phi)\bigr\rangle\Bigr\}.
\]

这意味着在 (10.2) 中我们可以在指数上插入一个常数 $1/\epsilon$ 而不改变答案。（无论如何，总有 $\Sigma$ 的体积这个因子进入作用量，从而进入 (10.2) 中被积函数的指数。）

% [Tutor Note: 「可以插入 $1/\epsilon$」的依据是 $L=\delta(\cdots)$ 的 $\delta$-恰当性：由于 $\delta$ 是作用量所对应的 BRST 型算子，$\int e^{-S-\frac1\epsilon\delta\Psi}=\int e^{-S}$（差一个拓扑不变量），取 $\Psi=\frac{1}{2\sqrt2}\langle\bar\chi',\operatorname{grad}W\rangle$ 便得到指数上的 $\frac1\epsilon\|\operatorname{grad}W\|^2$。注意这里的 $\bar\eta$ 与 $\delta$ 是同一套记号中的两个算子（$\bar\eta$ 为「上同调」方向、$\delta$ 为其伴随），不可与 (11.x) 中的超对称变分 $\delta$ 混为一谈。]

费米子积分很容易做出，于是我们把 $E$ 的计算归结为 $\mathbb{C}^n$ 上一个普通的积分：
\[
  E=\int_{\mathbb{C}^n}P(\phi)\,\epsilon^{-n}\det\operatorname{Hess}\overline{W}(\phi)\,
  e^{-\frac1\epsilon\|\operatorname{grad}W(\phi)\|^2}\frac{d^n\bar\phi\,d^n\phi}{(2\pi i)^n}.
\]

注意我们用 $\bar\omega$ 把 $\det\operatorname{Hess}\overline{W}(\phi)$ 解释为一个数。在 $\epsilon\to0$ 的极限下，剩下的积分局域化到 $W$ 的临界点集上。为简单起见，设 $W$ 具有孤立的非退化临界点。在临界点附近我们用 Gauss 积分近似该积分，这给出分母上的 $\|\det\operatorname{Hess}W(\phi)\|^2$。它与被积函数中的反全纯部分相消，最终结果是
\[
  E=\sum_{\text{临界点 }\phi_i}\frac{P(\phi_i)}{\det\operatorname{Hess}W(\phi_i)}.
\]

\subsection{习题 11 的解答}

在着手路径积分之前，我们先回顾\footnote{% 原文脚注 1
关于此计算的至少一部分，参考文献是 S.~K. Donaldson, P.~Kronheimer, \emph{The Geometry of Four-Manifolds}, Oxford University Press, New York, 1990, p.~197.}一个关于某种万有曲率（universal curvature）的公式。设 $P\to\Sigma$ 是 Riemann 流形 $\Sigma$ 上的主 $G$-丛。以 $\mathcal{A}_P$ 记 $P$ 上联络构成的空间，$\mathcal{G}_P$ 记规范变换群。以下我们自由地在联络 $A$ 与它所决定的 $\bar\partial$ 算子之间来回切换

% 原句跨页：p.301 开头「(complex structure) it defines on associated bundles by its (1,0) part」是本句的插入语，译文见下页首行。

% ===== 原文 p.301 =====
（即它由其 $(1,0)$ 部分在相伴丛上定义的复结构）。现在，丛 $\mathcal{A}_P\times_P X\to\mathcal{A}_P\times X$ 上有一个万有联络（tautological connection）。要把它下降到商\footnote{% 原文脚注 2
此项讨论中我们忽略奇点（singularities）。}$\mathcal{P}\to(\mathcal{A}_P/\mathcal{G}_P)\times X$ 上，我们需要 $\mathcal{A}_P\to\mathcal{A}_P/\mathcal{G}_P$ 上的一个水平分布，这可以通过取轨道（orbits）的正交补而得到。现在限制到 $\mathcal{M}_P\times X$，其中 $\mathcal{M}_P$ 是平坦联络模去规范变换所得的模空间。我们进一步限制到不可约联络。$\mathcal{M}_P$ 在 $A$ 的等价类处的切向量是上同调群 $H^1_A(\mathfrak{g}^P)$ 的一个元素，它由 $\Sigma$ 上一个取值于伴随丛的 1-形式 $\alpha$ 代表。给定两个这样的 $\alpha_1,\alpha_2$，此联络在 $\mathcal{P}\to\mathcal{M}_P\times X$ 上沿 $\mathcal{M}_P$ 的切方向的曲率 $\omega_1$ 为
\[
  \omega_1(\alpha_1,\alpha_2)=2\Delta_A^{-1}\bigl\langle[\alpha_1\wedge\alpha_2]\bigr\rangle,
\]
其中 $\Delta_A=d_A^*d_A$ 是 $\Omega^0(\mathfrak{g}^P)$ 上的 Laplace 算子（laplacian），而记号 $\langle[\alpha_1\wedge\alpha_2]\rangle$ 的意思是：对 $\alpha_i=a_i\otimes\zeta_i\in\Omega^1\otimes\mathfrak{g}$，
\begin{equation}
  \tag{11.1}
  \label{eq:sol11-1}
  \bigl\langle[a_1\otimes\zeta_1,a_2\otimes\zeta_2]\bigr\rangle=\langle a_1,a_2\rangle[\zeta_1,\zeta_2].
\end{equation}

若 $\Sigma$ 是 Riemann 面，我们用 $\square_A:\Omega^0(\mathfrak{g}^P)\to\Omega^2(\mathfrak{g}^P)$（由 $\square_A=-d_A*d_A$ 定义）把 (11.1) 改写为
\begin{equation}
  \tag{11.2}
  \label{eq:sol11-2}
  \omega_1(\alpha_1,\alpha_2)=2\square_A^{-1}[\alpha_1\wedge*\alpha_2].
\end{equation}

% [Tutor Note: 这里最容易出错的是算子的「次数」。$\Delta_A=d_A^*d_A$ 作用在 $\Omega^0$ 上（作为算子是负定的），故 $\Delta_A^{-1}$ 把 0-形式送回 0-形式；而 (11.2) 中的 $\square_A$ 被定义成 $\Omega^0\to\Omega^2$，它把 0-形式送到 2-形式（即 $\square_A f=\Delta f$ 配上体积元），因此 $\square_A^{-1}$ 只能吃 2-形式。$[\alpha_1\wedge*\alpha_2]$ 恰好是 2-形式（$\alpha_i$ 为 1-形式、$*\alpha_2$ 为 1-形式），所以 (11.2) 的两边都是标量函数。若把 $\square_A^{-1}$ 误当作通常的 Laplace 逆算子（$\Omega^0\to\Omega^0$），就会为形式次数不匹配而困惑，也容易丢掉 (11.2) 相对 (11.1) 中 Hodge 星号带来的一个符号。]

（尽管我们并不使用这一事实，但请注意 $\omega_1$ 只依赖于 $\Sigma$ 的共形结构。）

现在固定 Hermite 向量空间 $W$ 上的一个酉表示 $\rho:G\to U(W)$，并令 $W^P\to\Sigma$ 为相伴的 Hermite 向量丛。平凡（无穷秩）丛 $\mathcal{A}_P\times\Omega^0(W^P)$ 在限制到平坦联络之后，于规范变换下下降为丛 $\mathcal{S}\to\mathcal{M}_P$。这个丛也可以描述为 $W^P\to\mathcal{M}_P\times X$ 的相对截面（relative sections）丛。$\mathcal{P}\to\mathcal{M}_P\times X$ 上的联络诱导 $\mathcal{S}\to\mathcal{M}_P$ 上的一个联络：给定 $\mathcal{M}_P$ 的一个切向量，我们把它平凡地提升到 $\mathcal{M}_P\times X$，并逐点地在 $W^P$ 中作用协变导数。这个联络保持 $\mathcal{S}\to\mathcal{M}_P$ 上的 $L^2$ 度量，后者由 $W^P\to\Sigma$ 上逐点的内积对 $\Sigma$ 上的 Riemann 测度积分而得到。该联络的曲率由逐点地作用 $\dot\rho(\omega)$ 给出。我们实际上关心的是取值于 $W$ 的 $(0,1)$-形式，它是固定的复线丛 $\Lambda^{0,1}$ 与 $W^P$ 张量积的截面空间。联络与曲率都看不见这个线丛，所以公式 (11.2) 保持不变。现在考虑 $\Lambda^{0,1}\otimes W^P$ 的相对调和截面（relative harmonic sections）构成的子丛。它是一个有限维子丛\footnote{% 原文脚注 3
我们把讨论限制在使 $H^0(W^P)$ 为零的那些复结构上。}
\begin{equation}
  \tag{11.3}
  \label{eq:sol11-3}
  H^1(W^P)\to\mathcal{M}_P
\end{equation}
它是 $\mathcal{S}\to\mathcal{M}_P$ 的子丛，而我们可以用 $L^2$ 度量把联络投影到它上面。这个投影联络的曲率于是等于 (11.2) 逐点作用于全纯截面的投影，再加上一个第二基本形式（second fundamental form）项。以 $\pi_A$ 记到 $\Omega^{0,1}(W^P)$ 中关于联络 $A$ 调和的形式上的投影。稍作计算可知该第二基本形式项为
\begin{equation}
  \tag{11.4}
  \label{eq:sol11-4}
  \begin{aligned}
    \omega_2(\alpha_1,\alpha_2)
    &=\pi_A\,\dot\rho_{\mathbb{C}}\bigl(\alpha_2^{(0,1)}\bigr)\Delta_A^{-1}*\dot\rho_{\mathbb{C}}\bigl(\alpha_1^{(1,0)}\bigr)\pi_A\\
    &\quad-\pi_A\,\dot\rho_{\mathbb{C}}\bigl(\alpha_1^{(0,1)}\bigr)\Delta_A^{-1}*\dot\rho_{\mathbb{C}}\bigl(\alpha_2^{(1,0)}\bigr)\pi_A.
  \end{aligned}
\end{equation}

% [Tutor Note: (11.4) 的反对称性完全依赖于第二项前面的那个负号；$\omega_2$ 是 2-形式，必须满足 $\omega_2(\alpha_1,\alpha_2)=-\omega_2(\alpha_2,\alpha_1)$，读者可用此检验是否抄错了指标 $1\leftrightarrow2$。另注意算子的次序不可交换：$\dot\rho_{\mathbb{C}}(\alpha^{(0,1)})$ 先把 $(0,1)$-形式乘成取值于 $W$ 的形式，再经 $\Delta_A^{-1}*$ 变回 $(0,1)$-形式（星号把两个 $(0,1)$ 因子的乘积变回 $(0,1)$ 型），最后 $\pi_A$ 投到调和部分；若把 $\Delta_A^{-1}$ 与 $\dot\rho$ 交换，结果会多出与曲率有关的项。]

% ===== 原文 p.302 =====
(11.3) 的曲率为下述和
\begin{equation}
  \tag{11.5}
  \label{eq:sol11-5}
  \pi_A\omega_1\pi_A+\omega_2.
\end{equation}

现在转向物理。扭曲手征超场（twisted chiral superfield）有分量场（component fields）\footnote{% 原文脚注 4
在 [I-Supersolutions, §9.2] 中我们用的是 $\lambda_\pm$，它们与这里所用的 [I-Supersolutions, §6.2] 的 $\lambda^\pm$ 由
\[
  \lambda^+=-\lambda_-,\qquad \lambda^-=\lambda_+
\]
相联系。}$\sigma,\lambda^\pm,\bar\lambda^\pm$，它们取值于伴随丛 $\mathfrak{g}^P$。此外 $P$ 上还有一个联络 $A$。与习题 10 中一样，手征超场有分量场 $\phi,\psi_\pm,\bar\psi_\pm$，它们是相伴 Hermite 丛 $V^P$（或其共轭 $\overline{V}^P$）的截面。这里所用的拓扑扭曲与习题 10 中的不同；它与超对称 $\overline{Q}_++Q_-$ 对易。我们总结于下表：

\begin{table}[htbp]
  \centering
  \caption{扭曲手征超场各分量场的量子数。（原文表格，p.302）}
  \label{tab:p302}
  \renewcommand{\arraystretch}{1.18}
  \begin{tabular}{|c|c|c|c|c|}
    \hline
    分量场 & $K$ & $J_-$ & $J_+$ & $K+\frac12J_-+\frac12J_+$\\
    \hline\hline
    $A$ & $1\oplus-1$ & $0$ & $0$ & $1\oplus-1$\\
    \hline
    $\sigma$ & $0$ & $1$ & $-1$ & $0$\\
    \hline
    $\bar\sigma$ & $0$ & $-1$ & $1$ & $0$\\
    \hline
    $\lambda^+$ & $-1/2$ & $0$ & $-1$ & $-1$\\
    \hline
    $\bar\lambda^+$ & $-1/2$ & $0$ & $1$ & $0$\\
    \hline
    $\lambda^-$ & $1/2$ & $-1$ & $0$ & $0$\\
    \hline
    $\bar\lambda^-$ & $1/2$ & $1$ & $0$ & $1$\\
    \hline
    $\phi$ & $0$ & $0$ & $0$ & $0$\\
    \hline
    $\bar\phi$ & $0$ & $0$ & $0$ & $0$\\
    \hline
    $\psi_+$ & $1/2$ & $0$ & $1$ & $1$\\
    \hline
    $\bar\psi_+$ & $1/2$ & $0$ & $-1$ & $0$\\
    \hline
    $\psi_-$ & $-1/2$ & $1$ & $0$ & $0$\\
    \hline
    $\bar\psi_-$ & $-1/2$ & $-1$ & $0$ & $-1$\\
    \hline
  \end{tabular}
\end{table}

这里 $K$ 是旋转的无穷小生成元，$J_\pm$ 是 R 对称性（R-symmetries）。在曲面 $\Sigma$ 上的扭曲理论中，我们把 $\sigma,\bar\sigma$ 解释为取值于伴随丛的复标量；$\lambda^-,\bar\lambda^+$ 为奇的伴随值复标量；$\bar\lambda^-$ 为 $\Sigma$ 上取值于复化伴随丛的 $(1,0)$-形式；$\lambda^+$ 为 $\Sigma$ 上取值于复化伴随丛的 $(0,1)$-形式；$\phi$ 为 $V^P$ 的截面；$\bar\psi_+$（相应地 $\psi_-$）为 $\overline{V}^P$（相应地 $V^P$）的奇截面；$\psi_+$ 为 $\Sigma$ 上取值于 $V^P$ 的 $(1,0)$-形式；$\bar\psi_-$ 为 $\Sigma$ 上取值于 $\overline{V}^P$ 的 $(0,1)$-形式。

% [Tutor Note: 脚注 4 的等式在印刷本上左右两边的符号是同一套 $\lambda^\pm$，因此字面上是自我指涉的（$\lambda^+=-\lambda^-$ 与 $\lambda^-=\lambda^+$ 不能用同一组符号同时成立）。其原意应是：左边为 [I-Supersolutions, §9.2] 的下标记号 $\lambda_\pm$，右边为本讲所用、即 [I-Supersolutions, §6.2] 的上标记号 $\lambda^\pm$，即 $\lambda_+=-\lambda^-$、$\lambda_-=\lambda^+$。这一处记号冲突在对照该丛书第 I 卷使用时必须留意，否则 (11.6)–(11.7) 中的 $\pm$ 会整体反号。]

差 $J_--J_+$ 与 $\overline{Q}_++Q_-$ 对易，而用这个对称性来检查下面的 lagrangian 与方程式是否平衡是很方便的。

lagrangian 由四维中一般的 $\mathcal{N}=1$ lagrangian 经维数约化（dimensional reduction）得到。有关的公式出现在 [I-Supersolutions, §6.2] 中「$\mathbb{F}=\mathbb{C}$」的情形下。我们取那里的 lagrangian，维数约化到两维，然后 Wick 转动到欧氏空间。设

% 原句「Let ρ̇_C: g_C → End(V) denote ...」自 p.302 末行起，主语与谓语在下页首行。

% [Tutor Note: [I-Supersolutions] 指本丛书第 I 卷（\emph{Quantum Fields and Strings} 第 1 册）所附超对称习题解答一章；本书正文中凡 [I-...]、[II-...] 皆为此类卷内互引，页码随版本而异，核对时以节号（此处 §6.2、§9.2、§9.3）为准。]

% ===== 原文 p.303 =====
$\dot\rho_{\mathbb{C}}:\mathfrak{g}_{\mathbb{C}}\to\operatorname{End}(V)$ 为 $G$ 在 $V$ 上作用的复化无穷小作用（complexified infinitesimal action）。我们得到（维数约化后的 Yang–Mills 作用量见 [I-Supersolutions, §9.3]）：
\begin{equation}
  \tag{11.6}
  \label{eq:sol11-6}
  L_e=\left\{\begin{aligned}
    &\tfrac12|F_A|^2+\langle d_A\bar\sigma,d_A\sigma\rangle+\langle d_A\bar\phi,d_A\phi\rangle
      +\langle\bar\lambda\not D_A\lambda\rangle+\langle\bar\psi\not D_A\psi\rangle\\
    &+\tfrac12\|[\sigma,\bar\sigma]\|^2+2\|\dot\rho_{\mathbb{C}}(\sigma)\phi\|^2+2|\mu(\phi)|^2\\
    &-\sqrt2\bigl(\langle\bar\lambda^-,[\bar\sigma,\lambda^+]\rangle
      +\langle\bar\lambda^+,[\sigma,\lambda^-]\rangle\bigr)\\
    &+\sqrt2\bigl(\langle\bar\psi_-,\dot\rho_{\mathbb{C}}(\sigma)\psi_+\rangle
      +\langle\bar\psi_+,\dot\rho_{\mathbb{C}}(\bar\sigma)\psi_-\rangle\bigr)\\
    &+\sqrt2\bigl(\langle\dot\rho_{\mathbb{C}}(\lambda^a)^*\phi^*,\psi_a\rangle
      -\langle\bar\psi_{\dot a},\dot\rho_{\mathbb{C}}(\bar\lambda^{\dot a})\phi\rangle\bigr)
  \end{aligned}\right\}|d^2x|.
\end{equation}

% [Tutor Note: 最后一行是全式最易抄错的地方：求和指标在第一个括号里写成上标 $a$（$\lambda^a$、$\psi_a$ 一上一下缩并），在第二个括号里却写成带点的 $\dot a$（$\bar\lambda^{\dot a}$、$\bar\psi_{\dot a}$）。点指标与无点指标分别对应 $\pm$ 两个分量，且二者只能成对缩并、不能互相替换；若把 $\bar\psi_{\dot a}$ 写成 $\bar\psi_a$，就会在费米子符号上出错。另注意 $(\dot\rho_{\mathbb{C}}(\lambda^a))^*$ 是对算子的 Hermite 共轭，作者把星号写在括号外，不要误读成对 $\lambda^a$ 取共轭。]

在第二行中，$\mu:V\to\mathfrak{g}^*$ 是 $G$ 在 $V$ 上作用的矩映射（moment map）。最后一行的求和是对 $\pm$ 求和。倒数第二项的记号在 [I-Supersolutions, §6.2] 中有解释；它其实就是最后一项的复共轭。第三、四行的项来自把四维中的费米子动能项作维数约化；第二行中的第二项来自把四维中的 $\langle d_A\bar\phi,d_A\phi\rangle$ 作维数约化。类似地，我们把各场的超对称变换律维数约化到两维，作 Wick 转动，并特殊化到扭曲理论的超对称 $\delta=\eta(\overline{Q}_++Q_-)$。所得方程为（纯规范理论的变换律见 [I-Supersolutions, §9.2]）：
\begin{equation}
  \tag{11.7}
  \label{eq:sol11-7}
  \begin{array}{r@{\qquad}l}
    (\delta A)_{1,0}=2\eta\bar\lambda^- & \delta\lambda^+=-2\sqrt2\,\eta\,\bar\partial_A\sigma\\[3pt]
    (\delta A)_{0,1}=2\eta\lambda^+ & \delta\bar\lambda^+=\eta i\bigl(\tilde f_A+2\tilde\mu(\phi)-i[\sigma,\bar\sigma]\bigr)\\[3pt]
    \delta\sigma=0 & \delta\lambda^-=-\eta i\bigl(\tilde f_A+2\tilde\mu(\phi)\bigr)+i[\sigma,\bar\sigma])\\[3pt]
    \delta\bar\sigma=-\sqrt2\,\eta\,(\lambda^-+\bar\lambda^+) & \delta\bar\lambda^-=2\sqrt2\,\eta\,\partial_A\sigma\\[3pt]
    \delta\phi=-\sqrt2\,\eta\,\psi_- & \delta\psi_+=-2\sqrt2\,\eta\,\partial_A\phi\\[3pt]
    \delta\bar\phi=-\sqrt2\,\eta\,\bar\psi_+ & \delta\bar\psi_+=2\eta\,\dot\rho_{\mathbb{C}}(\bar\sigma)\bar\phi\\[3pt]
    & \delta\psi_-=2\eta\,\dot\rho_{\mathbb{C}}(\sigma)\phi\\[3pt]
    & \delta\bar\psi_-=-2\sqrt2\,\eta\,\bar\partial_A\bar\phi
  \end{array}
\end{equation}

% [Tutor Note: 原文 $\delta\lambda^-$ 一行右端的圆括号不配对：印刷为 $-\eta i(\tilde f_A+2\tilde\mu(\phi))+i[\sigma,\bar\sigma])$，末尾多出一个右括号。译文照原样保留以便对照。按上一行 $\delta\bar\lambda^+$ 的结构推断，原意应为 $-\eta i\bigl(\tilde f_A+2\tilde\mu(\phi)+i[\sigma,\bar\sigma]\bigr)$（即方括号项也在括号内），但这属于推测，未擅改。]

这里 $\tilde\mu:V\to\mathfrak{g}$ 是由矩映射 $\mu:V\to\mathfrak{g}^*$ 借助 $\mathfrak{g}$ 上的度量构造出来的。我们用到了标准算子
\[
  \partial=\frac12\Bigl(\frac{\partial}{\partial\tau}-i\frac{\partial}{\partial x}\Bigr)
\]
及其共轭（在 $\mathbb{R}^2$ 上）。$A$ 的曲率是\footnote{% 原文脚注 5
在 [I-Supersolutions, §9.2] 中我们把 Minkowski 空间 $M^2$ 中规范场 $A$ 的曲率写成
\[
  F_A=-f_A\,dt\wedge dx
\]
其中 $t,x$ 是 $M^2$ 上的标准坐标。在联络与曲率中都作代换 $\tau=it$，我们得到 (11.8)，其中 $\tilde f_A=if_A$。}
\begin{equation}
  \tag{11.8}
  \label{eq:sol11-8}
  F_A=\tilde f_A\,d\tau\wedge dx.
\end{equation}

% [Tutor Note: 这个 $\partial=\frac12(\partial_\tau-i\partial_x)$ 的约定决定了 $(1,0)$ 与 $(0,1)$ 的归属：它等于复坐标 $z=\tau+ix$ 下的 $\partial_z$，因此 $\partial$ 把标量变成 $(1,0)$-形式、$\bar\partial$ 变成 $(0,1)$-形式。用 (11.7) 校验：$\delta\psi_+=-2\sqrt2\eta\,\partial_A\phi$ 而 $\psi_+$ 是 $(1,0)$-形式（表中 $K=1/2$），$\delta\lambda^+=-2\sqrt2\eta\,\bar\partial_A\sigma$ 而 $\lambda^+$ 是 $(0,1)$-形式（$K=-1/2$），两者自洽。若改用常见的 $z=x+i\tau$，则 $\partial$ 与 $\bar\partial$ 互换，$K$ 的分配会整体反号，这是核对本页时最常见的错误来源。]

% ===== 原文 p.304 =====
在扭曲之后，这些公式在 Riemann 面 $\Sigma$ 上继续成立。

按照 §14.3.1 中所述的不动点思想，拓扑关联函数的泛函积分局域化到在超对称下不变的玻色子位形上。因此玻色子方程由令 $\delta(\text{费米子})$ 变换律的右端为零而得到。在我们的问题中，由 (11.7) 得到
\begin{equation}
  \tag{11.9}
  \label{eq:sol11-9}
  \begin{array}{c}
    \tilde f_A=-2\tilde\mu(\phi)\\[3pt]
    \partial_A\sigma=\bar\partial_A\sigma=0\\[3pt]
    \bar\partial_A\bar\phi=0\\[3pt]
    [\sigma,\bar\sigma]=0\\[3pt]
    \dot\rho_{\mathbb{C}}(\sigma)\phi=0.
  \end{array}
\end{equation}

现在，酉群 $U(V)$ 在 $V$ 上作用的矩映射由下述公式给出
\begin{equation}
  \tag{11.10}
  \label{eq:sol11-10}
  \mu(\phi)(T)=-\frac{i}{2}\langle\bar\phi,T\phi\rangle,\qquad \phi\in V,\quad T\in\mathfrak{u}(V).
\end{equation}

我们借助 $\dot\rho:\mathfrak{g}\to\mathfrak{u}(V)$ 作拉回，得到 $G$ 在 $V$ 上作用的矩映射。若 $\mathfrak{g}$ 没有可交换的因子，则这是唯一可能。此时，若表示是不可约的，则 $\mu(\phi)=0$ 的唯一解是 $\phi=0$。于是泛函积分被约化到 (11.3) 的零截面。我们把自己限制在这一情形。我们还进一步假设丛的选取使得 $H^0_A(W^P)=0$，其中
\[
  W=\overline{V},
\]
或者更确切地说，我们限制到使这一点成立的那些复结构。于是由 (11.9) 中第三个方程可得 $\bar\phi=0$。第一个方程随即蕴含 $A$ 是平坦的，因为由 (11.10) 矩映射为零。我们假设 $A$ 不可约，从而第二个方程蕴含 $\sigma=0$。其余方程不带来任何新的东西。总结起来，玻色不动点的方程为：
\[
  A\ \text{平坦}
\]
\[
  \phi=\sigma=0.
\]

现在来看费米子。按照 §14.3 中铺陈的一般原理，我们归结为对费米子的最低模积分。§14.4 中已说明，若 $A$ 不可约（我们正作此假设），则可以令复标量为零
\[
  \lambda^-=\bar\lambda^+=0.
\]
此外，在 $H^0_A(W^P)=0$ 的假设下我们有
\[
  \psi_-=\bar\psi_+=0.
\]
非零的费米子零模是
\[
  \chi:=\psi_+,\qquad \bar\chi:=\bar\psi_-,\qquad \lambda^+,\qquad \bar\lambda^-,
\]

% 原句以逗号结束于 p.304 末，其谓语「restricted to be harmonic 1-forms」在 p.305 首行。

% ===== 原文 p.305 =====
它们被限制为调和 1-形式。

我们现在可以计算一个拓扑可观测量的期望值了。如问题陈述中所述，并且从 (11.7) 显而易见，这些可观测量是 $\sigma$ 及其后代（descendents）的多项式，正如 §14.4 中讨论的纯规范理论那样。按拓扑场论的一般原理（§14.3.1），路径积分约化到由满足 (11.9) 的玻色位形与费米子零模所构成的超流形的一个邻域。玻色子 $A$ 与零模 $\lambda^+,\bar\lambda^-$ 一起构成 $\Pi T\mathcal{M}_P$，而如 §14.4 中所解释，$\Pi T\mathcal{M}_P$ 上的积分等价于 $\mathcal{M}_P$ 上一个微分形式的积分。对 $\chi,\bar\chi$ 的零模的积分是向量空间 $\Pi H^1_A(W^P)$ 上的费米子积分。我们把这些场的非零模、以及其余各场——$\phi,\bar\phi,\sigma,\bar\sigma,A,\bar\lambda^+,\lambda^-,\bar\psi_+,\psi_-$——在作用量的二次近似下沿到 $\Pi(T\mathcal{M}_P\times H^1_A(W^P))$ 的法丛积掉，从而化归为这些有限维积分。该积分是 Gauss 型的，因此可以显式做出。然后我们再计算 $\Pi H^1_A(W^P)$ 上的积分，从而化为 $\mathcal{M}_P$ 上一个微分形式的积分。只有在最后一步，拓扑可观测量才进入，因此我们的计算结果将是一个微分形式，它乘在对应于该可观测量的形式上。我们将认出这个形式就是 (11.3) 的 Euler 形式（Euler form）。

% [Tutor Note: $\Pi$ 是「奇（parity-reversed）化」的记号：$\Pi T\mathcal{M}_P$ 指把 $\mathcal{M}_P$ 的切丛的纤维改成奇（Grassmann）变量所得的超丛，故玻色子 $A$ 与费米子零模 $\lambda^+,\bar\lambda^-$ 可以同处一个超流形上；$\Pi H^1_A(W^P)$ 同理。$\Pi$ 不是投影算子，也不表示圆周，抄录时极易与 $\prod$ 或 $\pi$ 混淆。]

首先，我们做标量场 $\phi,\bar\phi,\sigma,\bar\sigma$ 上的 Gauss 积分。在 lagrangian (11.6) 中我们略去第二行的势能项，因为它们在这些标量场中是四次（quartic）的。作用量中剩下的是一个无穷维积分，它类似于下面 $\mathbb{C}^N$ 上的积分：
\[
  \iint_{\mathbb{C}^N}\frac{d\overline{\Phi}\,d\Phi}{(2\pi i)^N}
  e^{-(\overline{\Phi}Q\Phi+\overline{\Phi}L+\overline{L}\Phi)}
  =(\det Q)^{-1}e^{\overline{L}Q^{-1}L}.
\]

这里 $Q$ 是 Hermite 算子，故 $\overline{\Phi}Q\Phi$ 是实二次型；$L$ 是复线性算子，故 $\overline{\Phi}L+\overline{L}\Phi$ 是实线性型。在我们的问题中 $\Phi=\binom{\phi}{\sigma}$，并且
\[
  Q=\square_A
\]
\[
  L=\sqrt2\begin{pmatrix}
    \dot\rho_{\mathbb{C}}(\lambda^a)\psi_a\\[2pt]
    [\lambda^+,\bar\lambda^-]+B(\bar\chi,\chi)
  \end{pmatrix}
\]
其中双线性形式 $B$ 满足
\begin{equation}
  \tag{11.11}
  \label{eq:sol11-11}
  \langle\bar\chi,\dot\rho_{\mathbb{C}}(\bar\sigma)\chi\rangle
  =\langle\bar\sigma,B(\bar\chi,\chi)\rangle.
\end{equation}

注意 $L$ 在费米子中是二次的，所以 $\overline{L}Q^{-1}L$ 在费米子中是四次的。对 $A$ 的非零模也有类似的玻色积分。$A$ 的含费米子的相互作用项——$\langle\bar\lambda\not D_A\lambda\rangle$ 与 $\langle\bar\psi\not D_A\psi\rangle$——把 $\bar\lambda^\pm$ 与 $\lambda^\pm$ 配对、把 $\bar\psi_\pm$ 与 $\psi_\pm$ 配对。当我们限制到零模时，这样的配对为零，因此尽管这个积分给出有效作用量中费米子四次的项，那些项却不起作用。当然，$A$ 的非零模上的积分还给出一个逆行列式。

其次，我们对所有费米子的非零模积分。费米子的有效 lagrangian 形如
\[
  \langle\bar\lambda\not D_A\lambda\rangle+\langle\bar\psi\not D_A\psi\rangle
  +\overline{\mathcal{L}}\,\mathcal{Q}^{-1}\mathcal{L}
\]

% ===== 原文 p.306 =====
其中 $\mathcal{Q},\mathcal{L}$ 包含 $Q,L$ 以及对应于 $A$ 的相应形式。在 Gauss 近似下我们略去四次项 $\overline{\mathcal{L}}\mathcal{Q}^{-1}\mathcal{L}$，只剩下一个简单的 Gauss 积分。它给出分子上的行列式，恰好与玻色行列式相消。

我们剩下的是对零模 $\chi,\bar\chi,\lambda^+,\bar\lambda^-$ 的费米子积分，以及对 $\mathcal{M}_P$ 的一个玻色积分。如前所述，我们只做 $\chi,\bar\chi$ 上的积分，以约化到 $\Pi T\mathcal{M}_P$ 上的积分。有效 lagrangian 现在只涉及势项 $\overline{\mathcal{L}}\mathcal{Q}^{-1}\mathcal{L}$，而限制到零模上它是
\[
  -2\bigl\langle\dot\rho_{\mathbb{C}}(\bar\lambda^-)^*\bar\chi,\square_A^{-1}\dot\rho_{\mathbb{C}}(\lambda^+)\chi\bigr\rangle
  +2\bigl\langle B(\chi,\bar\chi),\square_A^{-1}[\lambda^+,\bar\lambda^-]\bigr\rangle.
\]

根据 (11.11)，我们可以把第二项改写为
\[
  -2\bigl\langle\bar\chi,\dot\rho_{\mathbb{C}}\bigl(\square_A^{-1}[\lambda^+,\bar\lambda^-]\bigr)\chi\bigr\rangle.
\]

（这是原作用量 (11.6) 中的一项，只是把 $\sigma$ 代换成了其运动方程的解 $\square_A\sigma=-\sqrt2[\lambda^+,\bar\lambda^-]$。）于是我们所拥有的积分是
\begin{equation}
  \tag{11.12}
  \label{eq:sol11-12}
  \iint_{H^1_A(W^P)}\frac{D\bar\chi\,D\chi}{(2\pi i)^m}
  e^{\int_\Sigma\langle\bar\chi,\omega\chi\rangle}
  =\frac{\det(\omega)}{(2\pi i)^m},
\end{equation}
其中 $m=\dim H^1(W)$，并且
\[
  -\omega=2\dot\rho_{\mathbb{C}}\bigl(\square_A^{-1}[\lambda^+,\bar\lambda^-]\bigr)
  +2\dot\rho_{\mathbb{C}}(\bar\lambda^-)\square_A^{-1}\dot\rho_{\mathbb{C}}(\lambda^+).
\]

在 (11.12) 中 $\omega$ 作用在 $\chi$ 上，而 $\chi$ 是 $(1,0)$ 型的。为与 (11.5) 比较，我们需要写出伴随算子 $\omega^*$，它作用在 $(0,1)$ 型的 $\bar\chi$ 上。回忆 $\dot\rho$ 在实向量上是斜 Hermite 的（skew-Hermitian），可见 $\omega$ 的伴随是
\[
  \omega^*=2\dot\rho_{\mathbb{C}}\bigl(\square_A^{-1}[\lambda^+,\bar\lambda^-]\bigr)
  -2\dot\rho_{\mathbb{C}}(\lambda^+)\square_A^{-1}\dot\rho_{\mathbb{C}}(\bar\lambda^-).
\]

由 (11.5)，对应于 $\omega^*$ 的微分形式是 $H^1(V^P)\to\mathcal{M}_P$ 的曲率的两倍，故 $-\omega^*$ 是共轭丛 $H^1(W^P)\to\mathcal{M}_P$ 的曲率形式的两倍。因此，(11.12) 等于 $H^1(W^P)\to\mathcal{M}_P$ 的曲率的行列式的 $(i/\pi)^m$ 倍，而这个行列式是该「障碍丛」（obstruction bundle）的 Euler 形式的 $2^m$ 倍\footnote{% 原文脚注 6
勤勉的读者也许能从更细致的计算中发现这个因子并不出现。想来这可以通过仔细处理 $\chi,\bar\chi$ 上积分的归一化而得到解决。}。

% [Tutor Note: 这里的 $2^m$ 是作者本人也承认的未定因子（见脚注 6）：Grassmann 测度 $D\bar\chi D\chi$ 的归一化与 Euler 形式的 $2^m$（Pfaffian）约定不匹配，故引用 (11.12) 时只应取「相差一个只依赖 $m$ 的常数」这层意思，不要把 $2^m$ 当作可验证的系数。另注意 (11.12) 左边 $\omega$ 吃的是 $(1,0)$ 型的 $\chi$，而要对接 (11.5) 必须先取伴随 $\omega^*$ 换成作用在 $(0,1)$ 型 $\bar\chi$ 上的算子；这一步隐含用 Hermite 度量把 $H^1$ 与其共轭等同，跳过伴随就直接比较会误用 $\dot\rho$ 的斜 Hermite 性。]
